\section{Introduction}
Scientific design requires solution fields for many choices of physical
parameters. Changing a boundary condition or diffusivity changes the temperature
throughout a domain. A surrogate model learns this map without repeating fine
simulations. Multifidelity learning adds cheaper coarse solutions
\citep{kennedy2000,howard2023,niu2024}, leaving a practical choice: which model
should use them? Transfer learning, joint training, and coarse field correction
use different fidelity relationships, and their accuracy can change as those
relationships weaken \citep{faza2026}.

We train a surrogate library and use additional known fine solutions to fit
ensemble weights. \method{} compares three
rules: select the most accurate individual model, average models according to
their individual errors, or fit averaging weights to minimize the combined
error. New parameters
produce a fine field without another coarse solve (Figure~\ref{fig:overview}).
Each dataset has its own fitted ensemble, with one weight per model fixed across
inputs and grid points. Automation selects and combines models within prescribed
training recipes.

An ensemble increases surrogate computation while reusing expensive field data.
For climate modeling, the Regional Aerosol Model Intercomparison Project
(RAMIP) reports indicative costs of approximately 100{,}000 CPU core hours for
CESM2.1 and 400{,}000 for UKESM1 per coupled simulation spanning roughly
36 simulated years \citep[Appendix~A]{wilcox2023ramip}. Core hours measure total
processor usage. These figures describe RAMIP simulations, not the production
cost of our ERA5 archive. \citet[Section~6]{niu2024} similarly describe their
extra inference overhead as ``small compared with the actual simulators''.
When generating fields dominates the budget, additional surrogate computation
can therefore be reasonable. \method{} reuses available training fields across
candidates and the same reserved fine examples across selection and combination
rules.

This approach builds on engineering surrogate ensembles
\citep{goel2007,acar2009,viana2009}, automated model libraries
\citep{feurer2015,erickson2020}, and comparative studies of multifidelity field
surrogates \citep{brunel2025}. Our contribution has three connected parts.
First, we construct a library of nine field surrogates by adapting established
architectures to transfer, pooled supervision, coarse field correction, and
fine only prediction. A common parameter to field interface makes these
different uses of training data available for selection and combination.
Second, we assemble a benchmark of 22 datasets from our simulations and imported
data, documenting their fidelity definitions and comparing models on common
evaluation cases within each dataset. Third, \method{} connects this library
and benchmark through an automated ensemble workflow with an explicitly counted
budget of additional fine solutions. The underlying architectural components
and ensemble weighting principles are established.

We first compare individual library models with additional baselines, then
ask whether ensemble fitting improves on choosing one model. Error analysis,
fitting loss controls, and ERA5 climate emulation \citep{hersbach2020,niu2024}
examine when combining predictions helps. This separates library accuracy from
reliable selection. Historical results informed development, so these
comparisons remain retrospective.
